\documentclass[11pt]{article}
\usepackage[margin=1in]{geometry}
\usepackage{amsmath,amssymb,amsthm}
\usepackage[hidelinks]{hyperref}

\newtheorem{theorem}{Theorem}
\newtheorem{lemma}{Lemma}

\title{A Literature-Implied Sharpening of the Complex Median Lemma in arXiv:2407.14988}
\author{fa\_banach\_001, agent\_lane\_06}
\date{2026-07-19}

\begin{document}
\maketitle

\paragraph{Status.}
\textbf{literature\_implied\_answer (complex median lemma sharpening; not main open problem).}
This note records a scoped identification: the complex median lemma in Feng--Wang--Zhou
\cite{FWZ2407} is a pushforward-measure instance of the planar orthogonal equipartition
or second pancake theorem. The identification improves the lemma's constant from
$1/16$ to the optimal $1/4$. It does not solve the paper's open noncommutative
paraproduct boundedness problem, nor the low-$p$ endpoint range left open for
Theorem 1.6 when $n=1$.

\paragraph{Source question/passage.}
Theorem 1.9 of \cite{FWZ2407} (TeX label \texttt{divideS}, proved in Section 8)
states that if $(\Omega,\mathcal F,\mu)$ is a measure space, $I\in\mathcal F$ has
finite measure, and $b:I\to\mathbb C$ is measurable, then there are two orthogonal
lines $L_1,L_2$ in $\mathbb C$ whose four closed quadrants $T_1,\ldots,T_4$ satisfy
\[
  \mu(\{x\in I:b(x)\in T_i\})\ge \frac{1}{16}\mu(I),\qquad i=1,\ldots,4.
\]
After the proof, the source remarks that proving the theorem first for simple
functions and then passing to general functions by a limit argument is not clear.

\paragraph{Supporting theorem.}
Fakhrutdinov--Musin \cite{FM2511} recall the planar second pancake theorem:
a finite-area planar mass admits two perpendicular straight lines cutting the
mass into four equal pieces. Their Theorem 1 proves the finite-point version
algorithmically. The measure form is the exact geometric statement needed after
pushing $\mu|_I$ forward under $b$.

\begin{theorem}[Optimal closed-quadrant complex median]
Let $(\Omega,\mathcal F,\mu)$ be a measure space, let $I\in\mathcal F$ have finite
measure, and let $b:I\to\mathbb C$ be measurable. Then there exist two perpendicular
lines $L_1,L_2\subset\mathbb C$ such that, if $Q_1,\ldots,Q_4$ are the four closed
quadrants determined by these lines, then
\[
  \mu(\{x\in I:b(x)\in Q_j\})\ge \frac{1}{4}\mu(I),
  \qquad j=1,\ldots,4.
\]
The constant $1/4$ is the best possible universal constant.
\end{theorem}

\begin{proof}
Let $M=\mu(I)$. If $M=0$ there is nothing to prove. Put
\[
  \nu=b_\#(\mu|_I),
\]
a finite Borel measure on $\mathbb C\simeq\mathbb R^2$ with total mass $M$.
It is enough to find the perpendicular cuts for $\nu$ and then pull them back
through $b$.

For compactly supported absolutely continuous planar measures, the second pancake
theorem gives perpendicular lines whose four quadrants have exactly mass $M/4$.
For a general finite Borel measure $\nu$, choose compactly supported absolutely
continuous measures $\nu_m$ with $\nu_m(\mathbb C)=M$ and $\nu_m\Rightarrow\nu$
weakly. Let $L_{1,m},L_{2,m}$ be perpendicular exact-quarter cuts for $\nu_m$.

The family of cuts has a convergent subsequence. Indeed, by tightness choose
a disk $B_R$ with $\nu_m(B_R)>M/2$ for all large $m$. Each of the two cutting
lines is a median line for $\nu_m$, since the two adjacent quadrants on either
side have total mass $M/2$. Hence neither line can miss $B_R$; otherwise one
of its halfplanes would have mass $<M/2$. Thus the offsets of both lines are
bounded, and since the two lines are perpendicular their intersection points
and directions lie in a compact parameter set. Passing to a subsequence, the
two lines converge to perpendicular lines $L_1,L_2$.

Let $Q_j$ be a closed quadrant for the limiting pair, and let $Q_{j,m}$ be the
corresponding closed quadrant for the approximating pair. For every $\delta>0$,
eventually $Q_{j,m}\subset (Q_j)_\delta$, the closed $\delta$-neighborhood of
$Q_j$. Therefore
\[
  \frac{M}{4}=\nu_m(Q_{j,m})\le \nu_m((Q_j)_\delta).
\]
Since $(Q_j)_\delta$ is closed and $\nu_m\Rightarrow\nu$, the Portmanteau theorem
gives
\[
  \frac{M}{4}\le \nu((Q_j)_\delta).
\]
Letting $\delta\downarrow0$ and using continuity from above for finite measures
yields $\nu(Q_j)\ge M/4$. Pulling back by $b$ gives the asserted inequality.

Optimality follows by applying any stronger putative constant to a line-null
planar mass, such as normalized area measure on a disk. The four open quadrants
are disjoint up to null boundary, so no universal lower bound can exceed $1/4$.
\end{proof}

\paragraph{Scope and novelty.}
The supporting authors do not present this as an answer to arXiv:2407.14988.
The implication is obtained by identifying the source's complex-valued median
construction with the classical orthogonal equipartition theorem for the
pushforward measure. The result is useful as a cleaner and sharper replacement
for Theorem 1.9, and it gives the source's desired limit argument in the
closed-quadrant formulation. It should not be counted as progress on the main
operator-valued open questions in the source paper.

\begin{thebibliography}{9}
\bibitem{FWZ2407}
X. Feng, T. Wang, and Y. Zhou,
\emph{On the boundedness and Schatten class property of noncommutative martingale
paraproducts and operator-valued commutators},
arXiv:2407.14988, 2024.
\url{https://arxiv.org/abs/2407.14988}

\bibitem{FM2511}
A. Fakhrutdinov and O. R. Musin,
\emph{Algorithms for orthogonal partitioning into four parts},
arXiv:2511.20866v2, 2026.
\url{https://arxiv.org/abs/2511.20866}
\end{thebibliography}

\end{document}
